Il progetto Prosody To Music nasce dall'idea di associare una composizione musicale a testo in lingua italiana.
Spesso, durante la lettura, può succedere di seguire un flusso melodico spontaneo, dando intonazione alle parole; difatti, la lingua italiana è di 
per sé fortemente musicale e ritmica.
Sfruttando dunque la prosodia ricavabile dal susseguirsi di parole, questo programma è in grado di produrre una base musicale 
coerente con l'input testuale.


\section{Obiettivi}
Il progetto ha come obiettivo principale la definizione di una pipeline coerente per trasformare un input testuale in una rappresentazione musicale espressiva 
e controllabile.
Nello specifico, il progetto mira al raggiungimento dei seguenti obiettivi:

\begin{enumerate}
    \item Prendere in input un testo ed estrapolare, secondo le regole della lingua italiana, varie tipologie di accenti, pause respiratorie e 
    significati sintattici.
    \item Produrre uno schema ritmico musicale coerente con la distribuzione degli accenti ritmici e dei silenzi.
    \item Generare una melodia coerente con il ritmo ottenuto precedentemente, lasciando libertà di personalizzazione all'utente.
\end{enumerate}
\section{Metodologia in breve}
Al fine di raggiungere gli obiettivi presentati sarà seguita la seguente pipeline:
\begin{enumerate}
    \item \textbf{Analisi del testo:} l'analisi del testo sarà svolta parallelamente dal punto di vista semantico e da quello ritmico per completare un processo di analisi prosodica del testo.  La semantica servirà a generare un'atmosfera musicale coerente con il significato e il contesto emotivo del testo. La prosodia servirà ad individuare gli accenti ritmici e le pause per la produzione dello schema ritmico.
    \item \textbf{Produzione dello schema ritmico:} traduzione dello studio prosodico in una struttura ritmicas coerente con il testo di input.
    \item \textbf{Addestramento del modello:} per la generazione della melodia, in questa fase si addestrerà un modello Transformer su un dataset di brani musicali.
    \item \textbf{Generazione musicale:} sarà generato un file MIDI coerente con il ritmo e con la semantica del testo, attraverso il Transformer precedentemente allenato.
    \item \textbf{Sintesi del brano:} il file MIDI sarà infine sintetizzato e convertito in un file audio (WAVE).      
\end{enumerate}


\section{Specifiche del progetto}
Il linguaggio di programmazione adottato per la realizzazione di questo progetto è Python, in quanto dotato di numerose librerie estremamte utili 
per l'implementazione. Inoltre, per l'addestramento dei modelli sono previsti dataset reperibili online contenenti dati e strutturati e pertinenti.
I dataset e le librerie saranno presentate nel dettaglio del capitolo successivo.